\documentclass[a4paper,10pt]{article}
\usepackage[utf8]{inputenc}

\usepackage[english]{babel}
\usepackage{xcolor}
\usepackage[compact,small]{titlesec}
\usepackage{booktabs}
\usepackage{multirow}
\usepackage{amsfonts,amsmath,amssymb}
\usepackage{marginnote}
\usepackage[top=1.8cm, bottom=1.8cm, outer=1.8cm, inner=1.8cm, heightrounded, marginparwidth=2.5cm, marginparsep=0.5cm]{geometry}
\usepackage{enumitem}
\usepackage{dialogue}
\setlist{noitemsep,parsep=2pt}
\newcommand{\highlight}[1]{\textcolor{kuleuven}{#1}}
\usepackage{pythonhighlight}
\usepackage{cleveref}
\usepackage{graphicx}
\usepackage{float}
\usepackage{caption}
\newcommand{\nextyear}{\advance\year by 1 \the\year\advance\year by -1}
\newcommand{\thisyear}{\the\year}
\newcommand{\deadlineGroup}{November 7, \thisyear{} at 9:00 CET}

\newcommand{\ReplaceMe}[1]{{\color{blue}#1}}
\newcommand{\RemoveMe}[1]{{\color{purple}#1}}

\setlength{\parskip}{5pt}

%opening
\title{Evolutionary Algorithms: Group 14 report}
\author{Borry Thomas, Chakhovich Mikhail, Geneste Victor and Van Vyve Lowie}

\begin{document}
\fontfamily{ppl}
\selectfont{}

\maketitle




\section{A basic evolutionary algorithm} 

\subsection{Representation}

For the Travelling Salesperson Problem (TSP), a permutation representation (or its mathematical equivalent, cycle notation) is adopted, as it provides a natural and direct way to encode the visiting order of cities. Each candidate solution corresponds to a specific permutation of city indices. For instance, the sequence [0, 2, 1, 3] represents a tour starting at city 0, visiting cities 2, 1, and 3 in that order, and then returning to the starting point. This representation offers a clear and compact structure that facilitates the generation, modification, and evaluation of tours. Although an adjacency representation was also considered, it was deemed less intuitive and more complex to manipulate. In the Python implementation, each permutation is stored as a \textit{NumPy ndarray}, which enables efficient numerical operations.

\subsection{Initialization}

The initial population is generated by creating random permutations of the city nodes, ensuring that each candidate solution corresponds to a valid and complete tour. In the implementation, random cycles are generated using the \textit{numpy.random.permutation} function, which efficiently produces random permutations of the city indices. This approach guarantees a varied and representative initial population for the optimization process.

\subsection{Selection operators}

Several selection operators were considered for the evolutionary process, including rank-based selection and k-tournament selection, while fitness-proportional selection was deliberately avoided due to its tendency to inadequately distribute the candidate solutions. Tournament selection was implemented because it effectively balances the exploitation of high-quality solutions with the preservation of population diversity. In this method, a subset of individuals is randomly sampled from the population, and the individual with the best fitness (i.e., lowest tour cost) within this subset is selected as a parent for reproduction. Tournament selection maintains diversity by allowing weaker solutions a non-zero probability of selection, thereby reducing the risk of premature convergence while still favoring fitter individuals. The tournament size is a key parameter, typically chosen between 2 and 5, to balance selective pressure with population diversity.

\subsection{Mutation operator}

Several mutation operators were considered, including insertion, swap, inversion, and scramble mutations. Swap mutation was implemented due to its simplicity and ease of implementation, while still introducing sufficient randomness to facilitate exploration of the solution space. In this method, two cities in a tour are selected and their positions exchanged, producing a slightly altered but valid tour. The degree of variation can be controlled through the mutation probability, which was set relatively high to ensure adequate exploration of new solutions. While swap mutation is effective for moderate-length tours, alternative mutation strategies may be required for longer tours to maintain diversity and prevent premature convergence.

\subsection{Recombination operator}

Several recombination operators were considered for the TSP, including order crossover (OX), partially mapped crossover (PMX), and edge recombination crossover (ERX). PMX was implemented because it preserves valid permutations while effectively combining segments from both parents. It operates by selecting two random crossover points in the first parent and copying the segment between them into the offspring. The remaining positions are filled by systematically mapping elements from the second parent to maintain the permutation property, following a mapping chain whenever conflicts occur. PMX effectively combines beneficial traits from both parents, particularly when shared subsequences exist. With little overlap between parents, it still produces valid offspring with increased variability, promoting exploration of the solution space. The operator can be influenced by parameters such as the number and position of crossover points (normally two random ones). No self-adaptive mechanism was implemented.

\subsection{Elimination operators}
The ($\lambda + \mu$) elimination strategy was implemented, as it effectively maintains genetic diversity in the parent generation while allowing exploitation by the next generation. This operator has two parameters: $\lambda$ representing the population size and $\mu$ the number of offspring generated in each iteration. In this implementation, $\lambda = 500$ and $\mu = 250$ were chosen. During elimination, individuals are removed from the combined set of parents and offspring until the population is reduced back to the size $\lambda$ ensuring that high-quality solutions are retained while diversity is preserved.

\subsection{Stopping criterion}

The evolutionary algorithm terminates based on two conditions. The primary condition is a time limit, monitored continuously using the built-in reporter function, with the algorithm continuing until the remaining time reaches zero. A logical secondary condition is convergence of the average fitness, which occurs when the population exhibits minimal variation, indicating that further iterations are unlikely to yield significant improvements. By combining these criteria, the algorithm utilizes the full allocated time unless the population has effectively converged, ensuring efficient use of computational resources while avoiding unnecessary evaluations. Although termination based on a fixed maximum number of iterations is theoretically possible, it may not achieve optimal performance within a given time window.


\section{Implementation generation}

\subsection{Prompts to generate code}
The core idea during the code generation phase was to decompose the implementation into modular components and prompt the LLM to generate each component individually. Additionally, to improve correctness and maintainability of the generated code, step-by-step (chain-of-thought) prompting was employed over multiple iterations. This strategy resulted in more precise prompts, better controllability, and a clearer structure in the final codebase.
\\\\
The TSP-representation was defined prior to prompting. After that, each component of the evolutionary algorithm was generated  by providing the function signature (function name, input specification, and required output). Below is an example of the prompting process used for generating the selection operator:

\begin{dialogue}
  \speak{Prompt 1}
    I want to create a selection operator for a genetic algorithm for the Travelings Salesman problem. The representation of the problem is already defined: a (permutation) list that represents the order of the visited cities. Which are the options and what are advantages and disadvantages?
  \speak{Effect on LLM}
    This ensures that the LLM “understands” the problem setup and has sufficient context to design a suitable selection operator. At the same time, it provides an overview of candidate approaches and their trade-offs.
  \speak{Prompt 2}
    I want you to implement k-tournament selection for a population of individuals in a genetic algorithm for the Travelings Salesman problem. Each individual is represented as a permutation list of cities. The function should take as input a list of individuals and an integer k, representing the number of individuals to select for the tournament. The function should return the selected individual from the tournament. To guide you, I already defined an Individual class that has a fitness attribute. You can use this information, together with the function header below, to implement the k-tournament selection.
    \begin{python}
    def tournament_selection(population: List[Individual], k: int) -> Individual:
    \end{python}
\end{dialogue}

\noindent
This multi-step prompting pattern led to code that closely matched the design expectations. Even for more complex operators such as the recombination method (Partially Mapped Crossover, PMX), the same approach proved effective. After receiving the generated code, it was manually reviewed and integrated into the overall project. Minor edits were made when necessary, such as renaming variables for consistency and adding concise clarifying comments.
%For instance, the refined version of Prompt 2 for recombination is shown below:

%\begin{dialogue}
%  \speak{Prompt 2}
%    I want you to implement partially mapped crossover for a population of individuals in a genetic algorithm for the Travelings Salesman problem. Each individual is represented as a permutation list of cities. To guide you, I already defined an Individual class that has a tour attribute, representing the order. You can create a new individual as follows: Individual(problem, order). The TravelingSalesmanProblem class has a $.get\_num\_cities()$ function. You can use this information, together with the function header below, to implement the recombination operator.
%    \begin{python}
%    def recombination(problem: TravelingSalesmanProblem, parent1: Individual, parent2: Individual ) -> Individual:
%    \end{python}
%\end{dialogue}

%\noindent
%After receiving the generated code, we manually reviewed and integrated it into the overall project. Minor edits were made when necessary, such as renaming variables for consistency and adding concise clarifying comments.

\subsection{Prompts to fix errors}

While the generated code was generally functional, several minor corrections were necessary during integration into the codebase. The most notable issue was a misalignment with the predefined data structures: the Individual class internally stores the tour as a \texttt{np.ndarray}, but the LLM initially implemented the order as a standard Python list. This inconsistency caused type mismatches when combining operators and computing distances. To resolve this, a corrective follow-up prompt was provided, asking the LLM to update the previously generated code accordingly:

\begin{dialogue}
  \speak{Prompt}
    You implemented the partially mapped crossover operator correctly, without any errors. However, there is a small issue: the order of an individual is implemented as a normal python array, while it should be a numpy array (np.ndarray). Please fix this in the code you generated before.
\end{dialogue}

\noindent
After receiving the updated response, it was verified that all operators produced consistent data structures and a few additional adjustments were made, such as renaming variables for readability and refining comments. Overall, postprocessing primarily involved ensuring internal consistency rather than modifying core logic.

\subsection{Critical reflection}


Generative AI excels at brainstorming and generating ideas, particularly when domain knowledge is limited. It can help get started quickly and provide a foundational implementation to build upon. However, effective use requires a solid understanding of the problem and underlying concepts, in order to guide the AI and critically evaluate its output. In this sense, prompt engineering becomes as important as traditional coding skills. While LLMs can accelerate the initial coding phase, they may also increase the time needed for reviewing, debugging, and integrating larger blocks of code. Furthermore, the AI does not always self-correct, so careful verification remains necessary. Given the rapid progress of LLMs in recent years, it has however become essential to at least learn how to work with them and understand what to realistically expect from their outputs. Ultimately, the developer's expertise and the trade-off between development speed and code quality determine the most effective way to leverage generative AI in a project. 


\section{Numerical experiments}

\subsection{Chosen parameter values}

%\ReplaceMe{What parameters are there to choose in your basic evolutionary algorithm? How did you approach parameter selection for now? Did you experiment a bit with different values? What was the influence of this choice? Which are the parameters that you will use for the experiments below?}

A critical component of this study was the systematic selection of parameters through a grid search over the primary control parameters of the evolutionary algorithm, including population size, offspring size, mutation rate, crossover probability, and tournament size. Each parameter configuration was evaluated over five independent runs to account for stochastic variability. For each combination, the mean and standard deviation of the best fitness values were recorded, measured after 2000 generations, to assess both performance and robustness of the algorithm. While larger population sizes generally led
to better performance, they also increased computation time. Higher mutation rates helped maintain diversity in the population, but too high rates could disrupt convergence. Table~\ref{tab:parameters} summarizes the tested ranges and highlights the selected values that were used in all subsequent experiments.


%One systematic approach to selecting parameters is to perform a grid search over a predefined set of values for each parameter. Each combination is tested multiple times under different settings to mitigate the influence of particularly lucky or unlucky runs. For our experiments, we varied the population size, offspring size, mutation rate, crossover rate, and tournament size. The specific values tested are summarized in Table \ref{tab: parameters}.
%\\\\
%Each run was allowed 1000 iterations to reach a solution. For each parameter combination and problem size, we recorded the mean and variance of the best fitness values obtained across repetitions.

%\textbf{CHECK THIS AND ADD MORE CONCLUSIONS: We observed that larger population sizes generally led to better performance, but also increased computation time. Higher mutation rates helped maintain diversity in the population, but too high rates could disrupt convergence.}

%\noindent
%We settled on the following parameters for our experiments as shown in Table \ref{tab: parameters}.

% \begin{table}[H]
%     \centering
%     \begin{tabular}{lc}
%         \hline
%         Parameter & Values \\
%         \hline
%         Population size & 50, \textbf{100}, 200 \\
%         Offspring size & 50, \textbf{100}, 200 \\
%         Mutation rate & 0.01, \textbf{0.05}, 0.1, 0.2 \\
%         Crossover rate & 0.6, \textbf{0.8}, 1.0 \\
%         Tournament size & 3, \textbf{5}, 7 \\
%         \hline
%     \end{tabular}
%     \caption{Parameters considered for grid search. Chosen values highlighted in bold.}
%     \label{tab:parameters}
% \end{table}

\begin{table}[H]
    \centering
    \begin{tabular}{lc}
        \hline
        Parameter & Tested values \\
        \hline
        Population size & 50, \textbf{100}, 200 \\
        Offspring size & 50, \textbf{100}, 200 \\
        Mutation rate $\alpha_\text{max}$ & 0.01, 0.05, 0.1, \textbf{0.12}, 0.2 \\
        Crossover probability $p_c$ & 0.6, \textbf{0.8}, 1.0 \\
        Tournament size $k$ & 3, \textbf{5}, 7 \\
        \hline
    \end{tabular}
    \captionsetup{width=0.8\textwidth}
    \caption{Parameter ranges explored during the grid search. Bold entries indicate the default configuration used in the subsequent experiments.}
    \label{tab:parameters}
\end{table}

% \subsection{Preliminary results}

%\ReplaceMe{Run your algorithm on the smallest benchmark problem. Include a convergence graph, by plotting the mean and best objective values in function of the time. How long did your program run until convergence? How do you rate the performance of your algorithm (time, memory, speed of convergence, diversity of population, etc)? What was the best fitness value you found? Do you think this is the global optimum? When you solve this problem several times, how much variation do you observe in the best solution? Include an informative plot of this.}
% \noindent
% As expected, Figure \ref{fig:convergence} shows that both the mean and best fitness values improve over time, indicating that the evolutionary algorithm is effectively optimizing the TSP solution. The algorithm converged within approximately ... iterations, taking in total ... seconds. This demonstrates a reasonable speed of convergence for this problem size. The fitness value of the best solution found was ..., which is close to the known benchmark of .... However, we cannot guarantee that this is the global optimum due to the NP-hard nature of the TSP. 

%\begin{figure}[H]
%    \centering
%    \includegraphics[width=0.7\textwidth]{convergence_plot.png}
%    \caption{Convergence graph showing mean and best objective values over time.}
%    \label{fig:convergence}
%\end{figure}
% \noindent
% The aspect of diversity of the population throughout the run plays a crucial role. Figure \ref{fig:convergence} illustrates that diversity decreases over time as the population converges towards better solutions. However, it remains above a certain threshold, indicating that the mutation and selection mechanisms are effectively maintaining some level of diversity to avoid premature convergence.
% \\\\
% When running the algorithm multiple times, we observed some variation in the best solution due to the stochastic nature of evolutionary algorithms. The standard deviation of the best fitness values across runs was ..., indicating a moderate level of variability. This suggests that while the algorithm is effective, there is still room for improvement in terms of consistency and robustness. The different scores are shown in Figure \ref{fig:variation}.

%\begin{figure}[H]
%    \centering
%    \includegraphics[width=0.7\textwidth]{variation_barchart.png}
%    \caption{Variation chart showing best fitness values across multiple runs.}
%    \label{fig:variation}
%\end{figure}


\subsection{Preliminary results}

Two separate experiments were conducted to evaluate the performance of the evolutionary algorithm.

\noindent
\textbf{1. Convergence analysis.} 
We first ran the GA using the selected parameter configuration on the smallest TSP benchmark instance to analyze convergence behavior. Figure~\ref{fig:convergence} shows the evolution of the mean and best fitness values over generations. The algorithm exhibits a rapid improvement during the first 200 iterations, followed by slower convergence toward a stable minimum after approximately 500 generations. Each run took on average around 18 seconds.  

\noindent
During these runs, population diversity steadily decreased and often fell to very low levels. This caused premature convergence and early stagnation, visible as plateaus in the fitness curves in Figure~\ref{fig:convergence}. Preserving diversity remains a key challenge to further improve the algorithm’s performance.

\begin{figure}[H]
    \centering
    \includegraphics[width=0.7\textwidth]{convergence_plot.png}
    \caption{Convergence of the mean and best fitness values over generations.}
    \label{fig:convergence}
\end{figure}

\noindent
\textbf{2. Repeatability study.} 
To evaluate stochastic variability and robustness, the algorithm was executed 100 times under identical settings. Each run recorded the best fitness achieved. Figure~\ref{fig:variation} shows a boxplot of these best fitness values. The red line indicates the median, and the green circle marks the mean. Whiskers show the approximate range of non-outlier values, while points outside the whiskers are considered outliers. Most runs achieved best fitness between approximately 17.500 and 19.500, indicating moderate stochastic variation but generally consistent convergence behavior.

\noindent
\begin{minipage}{0.4\textwidth}
Overall, these experiments demonstrate that the selected genetic algorithm configuration is capable of generating feasible solutions with moderate convergence speed and generally repeatable performance. Nevertheless, there remains considerable scope for improvement in surpassing heuristic solutions and approaching the global optimum. Population diversity decreases rapidly, resulting in premature convergence and early stagnation in some runs. The substantial variation in best fitness across the 100 independent executions complicates the identification of clearly optimal parameter settings, indicating that further investigation into parameter tuning and diversity-preserving mechanisms could enhance both solution quality and robustness. Future work may explore more sophisticated mutation operators, alternative selection strategies, local optimization techniques, or hybrid approaches to achieve a better balance between exploration and exploitation within the population.
\end{minipage}
\hfill
\begin{minipage}{0.55\textwidth}
    \begin{figure}[H]
        \centering
        \includegraphics[width=0.75\textwidth]{variation_boxplot.png}
        \caption{Boxplot of the best fitness values across 100 independent runs. The box shows the interquartile range, the red line is the median, the green circle indicates the mean. Whiskers represent the approximate range of non-outlier values.}
        \label{fig:variation}
    \end{figure}
\end{minipage}



\end{document}
